\section{System Design for Citizen Science}

\label{sec:system}
% ══════════════════════════════════════════════════════════════════════════════

\subsection{Deployment Architecture}

TaxoNet is deployed as a REST API service behind a geographic routing
layer. When a citizen scientist submits a photograph:

\begin{enumerate}
  \item \textbf{Candidate filtering}: Geographic and temporal priors
    reduce the candidate species set to 50--200 plausible species.
  \item \textbf{Feature extraction}: The ViT backbone extracts visual
    features in approximately 45\,ms on an NVIDIA T4 GPU.
  \item \textbf{Few-shot matching}: Taxonomic prototypical matching
    against the candidate set runs in approximately 12\,ms.
  \item \textbf{Result presentation}: Top-5 predictions with confidence
    scores and diagnostic field marks are returned to the user.
\end{enumerate}

Total inference latency is under 100\,ms per image.

\subsection{Prototype Management}

Species prototypes are computed offline from curated reference images and
stored in a vector database indexed by taxonomy and geography. New
species can be added by computing and inserting their prototype
embeddings---no model retraining is required.

\subsection{Active Learning Loop}

Citizen science identifications that disagree with model predictions
are flagged for expert review. Confirmed identifications update
prototypes:
\begin{equation}
  \mathbf{c}_{\text{new}}^{(s)} = (1 - \eta) \mathbf{c}_{\text{old}}^{(s)}
  + \eta \cdot \phi(x_{\text{confirmed}})
\end{equation}
where $\eta = 0.01$ is the update rate.


% ══════════════════════════════════════════════════════════════════════════════
